% Matches GraphHistoryActorCriticNet and Eqs. (encoders)--(policy).
\begingroup
\begin{tikzpicture}[
  x=1mm,y=1mm,>=Stealth,
  font=\fontsize{8.2bp}{9.8bp}\selectfont,
  inframe/.style={draw=black!80,line width=.55pt,fill=blue!9!white,
    align=center,inner sep=1.2mm,outer sep=0pt},
  modbox/.style={draw=black!80,line width=.55pt,fill=black!2!white,
    align=center,inner sep=1.2mm,outer sep=0pt},
  ctxbox/.style={draw=blue!55!black,line width=.5pt,fill=blue!3!white,
    align=center,inner sep=1.1mm,outer sep=0pt},
  outbox/.style={draw=black!80,line width=.55pt,fill=white,
    align=center,inner sep=1.2mm,outer sep=0pt},
  headfill/.style={fill=blue!7!white},
  dataedge/.style={-{Stealth[length=1.5mm,width=1.05mm]},
    draw=black!85,line width=.6pt},
  datawire/.style={draw=black!85,line width=.6pt},
  ctxedge/.style={-{Stealth[length=1.5mm,width=1.05mm]},
    draw=blue!55!black,line width=.55pt},
  ctxwire/.style={draw=blue!55!black,line width=.55pt},
  maskedge/.style={-{Stealth[length=1.5mm,width=1.05mm]},
    draw=black!65,line width=.55pt,dashed,dash pattern=on 1.3mm off .8mm},
  halo/.style={preaction={draw=white,line width=2.2pt}},
  lbl/.style={font=\fontsize{7.8bp}{9bp}\selectfont,inner sep=.8pt}
]
\path[use as bounding box] (0,0) rectangle (148,-112);

% 第1层：外部输入（候选特征与全局状态相互独立）
\node[inframe,anchor=north west,minimum width=108mm,minimum height=7mm]
  (enx) at (0,0)
  {\textbf{候选特征输入}\ $X_t\in\mathbb R^{M\times 41}$};
\node[inframe,anchor=north west,minimum width=35mm,minimum height=7mm]
  (engin) at (113,0)
  {\textbf{全局状态}\ $\bm g_t\in\mathbb R^{11}$};

% 第2层：特征编码（将6路属性合并为一个分组编码模块）
\node[modbox,anchor=north west,minimum width=64mm,minimum height=14mm]
  (enattr) at (0,-12.5)
  {\textbf{分组属性编码（6路）}\\
   发送方 $\bm v_m^s$、接收方 $\bm v_m^r$、语义模式 $\bm v_m^c$\\
   链路代价 $\bm v_m^l$、三方信念 $\bm v_m^b$、可靠性 $\bm v_m^q$};
\node[modbox,anchor=north west,minimum width=39mm,minimum height=14mm]
  (enbase) at (69,-12.5)
  {\textbf{基础编码}\\
   候选基础表示 $\bm h_m^0$};
\node[ctxbox,anchor=north west,minimum width=35mm,minimum height=14mm]
  (englob) at (113,-12.5)
  {\textbf{全局编码} $f_g$\\
   全局表示 $\bm z_t$};

\draw[dataedge] (32,-7) -- (32,-12.5);
\draw[dataedge] (88.5,-7) -- (88.5,-12.5);
\draw[ctxedge]  (130.5,-7) -- (130.5,-12.5);

% 第3层：上下文池化（仅对基础表示 h_m^0 加权池化）
\node[ctxbox,anchor=north west,minimum width=29mm,minimum height=11mm]
  (ensel) at (64,-32.5)
  {已选历史 $\bar{\bm h}_t^{\mathrm{sel}}$\\
   按 $e_m s_m^t$ 池化};
\node[ctxbox,anchor=north west,minimum width=29mm,minimum height=11mm]
  (enctr) at (96,-32.5)
  {中心上下文 $\bar{\bm h}_t^{\mathrm{ctr}}$\\
   按 $e_m c_m$ 池化};

% 基础编码输出分流：两路进上下文池化，一路向下进关系编码与门控残差
\draw[datawire] (88.5,-26.5) -- (88.5,-29.5);
\draw[datawire] (56,-29.5) -- (110.5,-29.5);
\fill[black!85] (88.5,-29.5) circle (.4mm);
\fill[black!85] (78.5,-29.5) circle (.4mm);
\draw[dataedge] (78.5,-29.5) -- (78.5,-32.5);
\draw[dataedge] (110.5,-29.5) -- (110.5,-32.5);

% 全局编码输出分流：一路进右侧共享上下文总线，一路向下进门控残差
\draw[ctxwire] (130.5,-26.5) -- (130.5,-29.5) -- (143,-29.5);
\fill[blue!55!black] (130.5,-29.5) circle (.4mm);

% 共享上下文总线（汇聚已选历史、中心上下文与全局表示）
\draw[ctxwire] (78.5,-43.5) -- (78.5,-46.5);
\draw[ctxwire] (110.5,-43.5) -- (110.5,-46.5);
\draw[ctxwire] (143,-29.5) -- (143,-73.5);
\draw[ctxwire] (50,-46.5) -- (143,-46.5);
\fill[blue!55!black] (78.5,-46.5) circle (.4mm);
\fill[blue!55!black] (110.5,-46.5) circle (.4mm);
\fill[blue!55!black] (143,-46.5) circle (.4mm);
\draw[ctxedge] (50,-46.5) -- (50,-53.5);

% 第4层：关系编码与门控残差融合
\node[modbox,anchor=north west,minimum width=60mm,minimum height=12mm]
  (enrel) at (0,-53.5)
  {\textbf{关系编码} $f_{\mathrm{rel}}$\\
   融合基础、属性与上下文，输出 $\bm u_m$};
\node[modbox,anchor=north west,minimum width=64mm,minimum height=12mm]
  (engate) at (72,-53.5)
  {\textbf{门控残差融合}\\
   $\bm h_m=\bm h_m^0+\bm a_m\odot\bm u_m$};

% 6路属性直连关系编码
\draw[dataedge] (22,-26.5) -- (22,-53.5);
\node[lbl,anchor=west] at (23.2,-38) {$\bm v_m^*$};

% h_m^0 连入关系编码，并作为残差跳连进入门控融合
\draw[datawire,halo] (56,-29.5) -- (56,-50);
\draw[datawire] (36,-50) -- (88,-50);
\fill[black!85] (56,-50) circle (.4mm);
\draw[dataedge] (36,-50) -- (36,-53.5);
\draw[dataedge] (88,-50) -- (88,-53.5);
\node[lbl,anchor=east] at (54.8,-38) {$\bm h_m^0$};

% z_t 连入门控残差融合（控制门控系数 a_m）
\draw[ctxedge,halo] (130.5,-29.5) -- (130.5,-53.5);
\node[lbl,anchor=west,text=blue!55!black] at (131.7,-38) {$\bm z_t$};

% 关系编码 -> 门控残差融合
\draw[dataedge] (60,-59.5) -- (72,-59.5);
\node[lbl,anchor=south] at (66,-58.8) {$\bm u_m$};

% 第5层：策略网络与价值网络
\fill[headfill] (0,-77) rectangle (64,-82.2);
\draw[black!60,line width=.4pt] (0,-82.2) -- (64,-82.2);
\draw[black!80,line width=.55pt] (0,-77) rectangle (64,-94);
\node[align=center,inner sep=0pt] at (32,-79.6) {\textbf{策略网络}};
\node[align=center,text width=61mm,inner sep=0pt] at (32,-88.1)
  {消息头：由 $\bm h_m$ 与上下文输出 $\ell_m$\\
   停止头：由上下文输出 $\ell_0$};

\fill[headfill] (74,-77) rectangle (136,-82.2);
\draw[black!60,line width=.4pt] (74,-82.2) -- (136,-82.2);
\draw[black!80,line width=.55pt] (74,-77) rectangle (136,-94);
\node[align=center,inner sep=0pt] at (105,-79.6) {\textbf{价值网络}};
\node[align=center,text width=59mm,inner sep=0pt] at (105,-88.1)
  {对合法候选 $\bm h_m$ 求掩码均值\\
   拼接上下文后输出价值 $V_\psi(s_t)$};

% 共享上下文总线接入策略网络与价值网络
\draw[ctxwire] (143,-73.5) -- (46,-73.5);
\fill[blue!55!black] (120,-73.5) circle (.4mm);
\draw[ctxedge] (120,-73.5) -- (120,-77);
\draw[ctxedge] (46,-73.5) -- (46,-77);
\node[lbl,fill=white,text=blue!55!black] at (69,-73.5)
  {上下文 $(\bar{\bm h}_t^{\mathrm{sel}},\bar{\bm h}_t^{\mathrm{ctr}},\bm z_t)$};

% 门控输出 h_m 接入策略网络与价值网络
\draw[datawire] (92,-65.5) -- (92,-69) -- (22,-69);
\fill[black!85] (92,-69) circle (.4mm);
\draw[dataedge] (22,-69) -- (22,-77);
\draw[dataedge,halo] (92,-69) -- (92,-77);
\node[lbl,fill=white] at (34,-69) {$\bm h_m$};

% 第6层：动作概率输出与合法动作掩码
\node[outbox,anchor=north west,minimum width=64mm,minimum height=11mm]
  (enpol) at (0,-100)
  {屏蔽非法消息后归一化为动作概率\\
   $\pi_\theta(a\mid s_t),\quad a\in\{0,1,\ldots,M\}$};
\node[inframe,anchor=north west,minimum width=54mm,minimum height=11mm]
  (enmask) at (78,-100)
  {\textbf{合法动作掩码}\ $\bm\mu_t$\\
   屏蔽重复与超预算项，停止动作恒可用};

\draw[dataedge] (32,-94) -- (32,-100);
\draw[maskedge] (78,-105.5) -- (64,-105.5);
\draw[maskedge] (105,-100) -- (105,-94);
\end{tikzpicture}
\endgroup